\subsection{Formally smooth quotients of algebras}

THEOREM 2.--- Let k be a ring, A a k-algebra, and J an ideal of A such that the k-algebra A/J is formally smooth. Endow A with the J-adic topology. The following conditions are equivalent:

\begin{enumerate}
\def\labelenumi{(\roman{enumi})}
\item
  the topological k-algebra A is formally smooth;
\item
  the \(A/\mathrm{J}\)-module \(\mathrm{J}/\mathrm{J}^2\) is projective and the canonical homomorphism (§ 5, no. 2)
\end{enumerate}

\[
\beta : \mathbf {S} _ {\mathrm{A/J}} (\mathrm{J} / \mathrm{J} ^ {2}) \rightarrow \operatorname{gr} _ {\mathrm{J}} (\mathrm{A})
\]

is bijective;

\begin{enumerate}
\def\labelenumi{(\roman{enumi})}
\setcounter{enumi}{2}
\tightlist
\item
  the A/J-module J/J \(^{2}\) is projective and there exists an isomorphism of topological k-algebras from the separated completion of A onto the completion of the graded algebra \(\mathsf{S}_{A/J}\) (J/J \(^{2}\)).
\end{enumerate}

If A is Noetherian, these conditions are also equivalent to:

\begin{enumerate}
\def\labelenumi{(\roman{enumi})}
\setcounter{enumi}{3}
\tightlist
\item
  the ideal J is completely secant.
\end{enumerate}

First observe that (iii) implies (i): indeed, under the hypotheses of (iii), the algebra \(\mathsf{S}_{\mathrm{A}/\mathrm{J}}(\mathrm{J}/\mathrm{J}^{2})\) , endowed with the topology associated with its grading, is formally smooth over A/J (no. 3, example 1), hence over k (no. 2, prop. 3, a)); assertion (i) then follows from prop. 3, c) of no. 2.

Let \(\widehat{\mathbf{A}}\) be the separated completion of the algebra A and \(\widehat{\mathbf{J}}\) the separated completion of J. The canonical homomorphism \(i: \mathbf{A} \to \widehat{\mathbf{A}}\) induces an isomorphism \(\mathrm{A} / \mathrm{J} \longrightarrow \widehat{\mathrm{A}} / \widehat{\mathrm{J}}\) (III, § 2, no. 12, formula (21)). Let \(\varphi: \mathrm{A} / \mathrm{J} \to \widehat{\mathrm{A}}\) be a lifting of this isomorphism (no. 4, prop. 5). Denote \(\lambda: \widehat{\mathrm{J}} \to \mathrm{J} / \mathrm{J}^2\) the surjection deduced from the canonical isomorphism \(\mathrm{J} / \mathrm{J}^2 \longrightarrow \widehat{\mathrm{J}} / \widehat{\mathrm{J}}^2\) (III, § 2, no. 12, formula (21)). Let \(a\) be an element of A, \(\bar{a}\) its class in \(\mathrm{A} / \mathrm{J}\), and let \(z\) be an element of \(\widehat{\mathrm{J}}\); one has \(\varphi(\bar{a}) \equiv i(a)\) (mod. \(\widehat{\mathrm{J}}\)), whence \(\varphi(\bar{a})z \equiv i(a)z\) (mod. \(\widehat{\mathrm{J}}^2\)) and \(\lambda(\varphi(\bar{a})z) = \lambda(i(a)z) = \bar{a}\lambda(z)\). In other words, \(\lambda\) is A/J-linear when one equips \(\widehat{\mathbf{J}}\) with the structure of an A/J-module induced by \(\varphi\).

Suppose that the homomorphism \(\lambda\) admits an A/J-linear section \(\sigma : J/J^{2} \longrightarrow \widehat{J}\). Let \(S\) be the \(k\)-graded algebra \(S_{A/J}(J/J^{2})\) and \(\widehat{S}\) its completion. Let

\[
\theta : \mathsf {S} \to \widehat {\mathsf {A}}
\]

the homomorphism of \(k\)-algebras such that \(\theta(x) = \varphi(x)\) for \(x\) in \(\mathbf{S}^0 = \mathrm{A}/\mathrm{J}\), and \(\theta(x) = \sigma(x)\) for \(x\) in \(\mathbf{S}^1 = \mathrm{J}/\mathrm{J}^2\). Since \(\theta\) maps \(\mathbf{S}^1\) into \(\widehat{\mathbf{J}}\), it maps \(\mathbf{S}^n\) into \(\widehat{\mathbf{J}}^n\) and therefore extends to a continuous homomorphism \(\widehat{\boldsymbol{\theta}}: \widehat{\mathbf{S}} \to \widehat{\mathbf{A}}\). The map \(\mathrm{gr}_1(\boldsymbol{\theta}): \mathrm{J}/\mathrm{J}^2 \longrightarrow \widehat{\mathrm{J}}/\widehat{\mathrm{J}}^2\) is the composite of \(\sigma\) with the canonical surjection \(\widehat{\mathrm{J}} \longrightarrow \widehat{\mathrm{J}}/\widehat{\mathrm{J}}^2\); since \(\sigma\) is a section of \(\lambda\), \(\mathrm{gr}_1(\boldsymbol{\theta})\) coincides with the canonical isomorphism from \(\mathrm{J}/\mathrm{J}^2\) onto \(\widehat{\mathrm{J}}/\widehat{\mathrm{J}}^2\). Consequently \(\mathrm{gr}(\boldsymbol{\theta}): \mathbf{S} \to \mathrm{gr}_{\widehat{\mathbf{J}}}(\widehat{\mathbf{A}})\) is the composite of the canonical surjection \(\beta\) with the canonical isomorphism \(\mathrm{gr}_\mathrm{J}(\mathrm{A}) \to \mathrm{gr}_{\widehat{\mathbf{J}}}(\widehat{\mathbf{A}})\) (III, § 2, no. 12, formula (22)).

Now let us prove the implication (ii)⇒(iii). Under hypothesis (ii), the A/J-module J/J² is projective, so λ admits an A/J-linear section; the homomorphism \(\widehat{\boldsymbol{\theta}}:\widehat{\mathbf{S}}\to\widehat{\mathbf{A}}\) associated with this section by the preceding construction induces, by hypothesis, an isomorphism on the associated graded modules, hence is bijective (III, § 2, no. 8, cor. 3 of th. 1), which proves (iii).

Let us prove (i)⇒(ii). Suppose the topological k-algebra A is formally smooth. First we prove that the A/J-module \(J/J^{2}\) is projective. Let M be an \(A/J\)-module and \(f:M\to J/J^{2}\) a surjective A/J-linear map; the goal is to prove that f admits an A/J-linear section.

Let \(\pi: A/J^{2} \to A/J\) be the canonical surjection. By Remark 1 of No. 2, there exists an isomorphism of \(k\)-algebras \(\psi: A/J \oplus J/J^{2} \longrightarrow A/J^{2}\) such that \(\pi(\psi(y,z)) = y\) and \(\psi(0,z) = z\) for \(y \in A/J\), \(z \in J/J^{2}\). Consider the \(k\)-algebra (A/J) \(\oplus M\) (No. 1, example) and the map \(u: (A/J) \oplus M \to A/J^{2}\) such that \(u(x,m) = \psi(x,f(m))\). This is a surjective homomorphism of \(k\)-algebras, whose kernel is the kernel of \(f\) in M, hence is nilpotent of square zero. The canonical surjection \(\rho: A \to A/J^{2}\) is continuous; since A is formally smooth as a topological \(k\)-algebra, there exists a homomorphism of \(k\)-algebras \(\tilde{\rho}: A \to (A/J) \oplus M\) such that \(u \circ \tilde{\rho} = \rho\). Since \(\mathrm{pr}_{1} = \pi \circ \psi = \pi \circ u\), one has \(\mathrm{pr}_{1} \circ \tilde{\rho} = \pi \circ u \circ \tilde{\rho} = \pi \circ \rho\), so that \(\mathrm{pr}_{1} \circ \tilde{\rho}\) is the canonical surjection of A onto A/J. We therefore have \(\tilde{\rho}(J) \subset M\) and consequently \(\tilde{\rho}(J^{2}) = 0\), so that \(\tilde{\rho}\) induces an A/J-linear map \(s: J/J^{2} \to M\). One has \(u \circ \tilde{\rho} = \rho\) and \(\mathrm{pr}_{2} \circ \psi^{-1} \circ u(y,m) = f(m)\) for \(y \in A/J\) and \(m \in M\). Let \(x \in J\), and let \(\bar{x}\) be its class in \(J/J^{2}\); one has \(f(s(\bar{x})) = f(\mathrm{pr}_{2}(\tilde{\rho}(x))) = \mathrm{pr}_{2}(\psi^{-1}(\bar{x})) = \bar{x}\). Thus \(s\) is a section of \(f\).

It remains to prove that the homomorphism \(\beta\) is injective. Since the A/J-module \(J/J^{2}\) is projective, \(\lambda\) admits an A/J-linear section; denote by \(\theta: S \to \widehat{A}\) the associated homomorphism. The homomorphism gr(θ) is identified with \(\beta\). Let \(m\) be an integer; denote by \(\Sigma_{m}\) the graded \(k\)-algebra quotient of \(S\) by the ideal \(\sum_{i > m} S^{i}\) and \(\theta_{m}: \Sigma_{m} \to A/J^{m+1}\) the homomorphism deduced from \(\theta\). The composite of \(\theta_{m}\) with the canonical surjection \(A/J^{m+1} \longrightarrow A/J\) is the canonical projection of \(\Sigma_{m}\) onto \(S^{0} = A/J\); consequently the kernel of \(\theta_{m}\) is a nilpotent ideal. By the example in no. 4, there exists a lifting \(\psi_{m}: A \to \Sigma_{m}\) of the canonical surjection \(A \to A/J^{m+1}\). Since the composite of \(\psi_{m}\) with the canonical projection from \(\Sigma_{m}\) onto A/J is the canonical surjection, \(\psi_{m}(J)\) consists of elements of degree \(>0\). By passing to the associated graded objects, we deduce from \(\psi_{m}\) a \(k\)-linear graded map gr(\(\psi_{m}\)): gr \(_{J}(A) \to \Sigma_{m}\) such that gr \(_{m}(\theta) \circ\) gr \(_{m}(\psi_{m}) = \text{Id}_{J^{m}/J^{m+1}}\). It follows that gr \(_{m}(\theta)\), hence also \(\beta_{m}\), is injective, which completes the proof of (ii).

Finally, when A is Noetherian, conditions (ii) and (iv) are equivalent (§ 5, no. 2, th. 1).

COROLLARY 1.--- Let \(k\) a field and A a \(k\) -local Noetherian algebra such that the extension \(\kappa_{A}\) of \(k\) Let it be separable. The following conditions are equivalent:

\begin{enumerate}
\def\labelenumi{(\roman{enumi})}
\item
  the \(k\) -algebra A is formally smooth for the topology \(\mathfrak{m}_{\mathrm{A}}\) -adic;
\item
  the ring A is regular;
\item
  the \(k\)-algebra A is absolutely regular (\(\S 6, \mathrm{n}^{\circ} 4\), def. 1);
\item
  the k-algebra \(\hat{A}\) is isomorphic to \(\kappa_{A}[[T_{1},\ldots,T_{n}]]\) , with \(n=\dim A\) .
\end{enumerate}

Conditions (ii) and (iii) are equivalent by Example 3 of § 6, No. 4, and amount to saying that the ideal \(\mathfrak{m}_{\mathrm{A}}\) is completely secant (VIII, § 5, no. 2, th. 1). Moreover, every ring isomorphism of \(\widehat{\mathrm{A}}\) on \(\kappa_{\mathrm{A}}[[\mathrm{T}_{1},\ldots,\mathrm{T}_{n}]]\) is bicontinuous, since these are local rings. Since the \(k\) -algebra \(\mathrm{A}/\mathfrak{m}_{\mathrm{A}}\) is formally smooth (no. 3, th. 1), the corollary follows from th. 2 applied with \(\mathrm{J}=\mathfrak{m}_{\mathrm{A}}\) .

COROLLARY 2.--- Let k be a field, A a Noetherian k-algebra, and J an ideal of A contained in the radical of A. Suppose the k-algebra A is formally smooth for the J-adic topology. Then it is absolutely regular.

Indeed, let \(k'\) be a finite extension of \(k\) and \(A'\) the A-algebra \(A_{(k')}\); it remains to prove that, for every maximal ideal \(\mathfrak{m}'\) of \(A'\), the Noetherian local ring \(A_{\mathfrak{m}'}'\) is regular. Now we have \(JA' \subset \mathfrak{m}'\): indeed, the inverse image of \(\mathfrak{m}'\) in A is a maximal ideal of A (V, § 2, no. 1, prop. 1), hence contains J. The \(k'\)-algebra \(A'\) is formally smooth for the \(JA'\)-adic topology (no. 2, prop. 4, b)), and the \(k'\)-algebra \(A_{\mathfrak{m}'}'\) is formally smooth for the \(JA_{\mathfrak{m}'}'\)-adic topology (no. 2, prop. 4, a)), hence also for the \(\mathfrak{m}'A_{\mathfrak{m}'}'\)-adic topology. Let \(k_0\) be the prime subfield of \(k'\). Then \(A_{\mathfrak{m}'}'\) is formally smooth over \(k_0\) for the \(\mathfrak{m}'A_{\mathfrak{m}'}'\)-adic topology (cor. of th. 1 of no. 3); since \(\kappa(\mathfrak{m}')\) is separable over \(k_0\), the ring \(A_{\mathfrak{m}'}'\) is regular (cor. 1).

COROLLARY 3.--- Let k be a ring and A a formally smooth k-algebra.

\begin{enumerate}
\def\labelenumi{\alph{enumi})}
\item
  The A-module \(\Omega_k(A)\) is projective.
\item
  Suppose that the ring A ⊗k A is Noetherian. Let μ : A ⊗k A → A be the homomorphism such that μ(x ⊗ y) = xy; then the ideal Ker(μ) is completely secant.
\end{enumerate}

The \(k\)-algebras A and \(\mathbf{A} \otimes_{k} \mathbf{A}\) are formally smooth (no. 2, prop. 4, c)), and A is isomorphic to the quotient of \(\mathbf{A} \otimes_{k} \mathbf{A}\) by the kernel I of \(\mu\). By definition, \(\Omega_{k}(\mathbf{A}) = \mathbf{I}/\mathbf{I}^{2}\). Thus a) and b) follow from th. 2.

